\documentclass[12pt,reqno]{amsart}
\usepackage{../handout, ../customcommands}
\usepackage{caption}

\hdr{MAT 240}{\S1.7 Counting methods}

\begin{document}

\textbf{Suggested problems:} 1-10

\bigskip
\hrule

\bigskip

\prob Suppose that we toss a coin six times. What is the cardinality of the sample space consisting of the outcomes of the six tosses?

\vfill
\prob A standard combination lock has a dial with tick marks for $40$ numbers from $0$ to $39$. The combination consists of a sequence of three numbers that must be dialed in the correct order to open the lock. What is the cardinality of the sample space consisting of all possible combinations?

\vfill
\prob Consider an experiment in which a card is selected and removed from a deck of $n$ different cards, a second card is then selected and removed from the remaining $n-1$ cards, and finally a third card is selected from the remaining $n-2$ cards. Each outcome consists of the three cards in the order selected. What is the cardinality of the sample space consisting of all possible selections of three cards?

\vfill
\newpage
\prob Suppose that a club consists of 25 members and that a president and a secretary are to be chosen from the membership. Determine the total number of ways in which these two positions can be filled.

\vfill
\prob (The Birthday Problem.) What is the smallest number $k$ such that, in a room of $k$ people, the probability that at least two people share a birthday is greater than $1/2$? 
\vfill


\end{document}